\chapter{Computation: Universality, Limits, and Stability}
\label{chap:computation-universality-limits-stability}

Here the apparatus runs closest to the machine. Run it as a computer and it
gains a single universal power and meets four kinds of wall: a wall of continuation where it must stop, a wall
of extraction where it cannot compress, a wall of decision where it cannot settle, and a wall of stability
where finite precision defeats it --- the program ours to write, the walls it runs into the world's. Eight
effects make the case, and one of them, the precision limit von
Neumann named, un-folds into two: the machine can replay anything it can write down, but it cannot always
continue, compress, decide, or stabilize.

\section{The Turing Effect: serialization and universality}
\label{se:turing}

Computation is universal because records serialize. Turing's insight was that a machine need not be built anew
for each task: one machine, reading a description off a tape, can imitate any other --- the universal machine,
the abstract ancestor of every stored-program computer. The structural fact beneath it is that any record,
whatever its dimensionality, can be written out as a sequence and later rebuilt from that sequence without
loss,
\begin{equation}
  \text{reconstruct}(\text{serialize}\, g) = g,
\end{equation}
the round-trip an identity --- nothing lost in flattening a structure to a line and rebuilding it --- the
flattening ours, what survives it the world's. A single
universal reader, taking one step per cell, can then replay any record so written. The universality of
sequential computation is just this: a record faithfully serialized and faithfully reconstructed loses nothing,
so one machine that reads sequences can run them all. The reach of that one machine is the content of the
Church--Turing thesis: every procedure anyone would recognize as a mechanical computation --- every effective
method, on any machine yet built or imagined --- can be carried out by the universal reader. Universality is not
the property of one clever construction but, if the thesis holds, the shape of computation itself.

The honest floor is name-only --- universality is labelled here, written as the physics of computation, with
the finite model claiming only the name --- the name ours to give, the reach the world's. The reading is that
universality is serialization, not magic. A record
reduces to a sequence and rebuilds from it without loss, and the universal machine is just the reader of those
sequences --- one step per cell, every record within its reach because every record can be written down.

\section{The Halt Effect: where refinement must stop}
\label{se:halt}

Not every refinement has a successor. There are configurations from which no admissible next step exists, and
there the refinement halts. The structure is a gate: before each step the order asks whether any successor may
follow --- whether some admissible event may precede the rest without contradicting what is already written ---
and where that gate returns nothing, the refinement stops --- the steps ours to take, the gate's answer the
world's. A halted ledger is not incomplete; it is complete in
the only sense the order allows, having reached a configuration with no admissible continuation. This inverts
the halting problem's usual lesson: there, that one cannot in general foretell whether a machine halts is
a limit on foresight; here, the halt itself is no failure but a terminus, the record finished because the order
has nowhere left to extend it.

The honest floor is a no-go --- when no successor satisfies the order, the process stops, and the stopping is
the order's wall, not the machine's act --- the running ours, the terminus the world's. The reading is that
halting is completion, not failure. The machine does not always run on;
where the order admits no next step it must stop, and that stopping is forced by the structure alone --- a
wall the refinement reaches and cannot pass.

\section{The Jupyter Effect: a collision halts the kernel}
\label{se:jupyter}

Asynchronous writes can collide, and a collision halts the kernel --- the kernel here the programmer's
everyday model, its name borrowed for the ledger's reader. When two updates target the same slot at the
same point in the order but carry different values --- same slot, same order, different values --- there is no
admissible reconciliation, and the kernel does not resolve the conflict but halts --- the writes ours to
issue, the refusal the world's. The future cannot be
accessed before the past, and a contradiction at one slot has no successor that respects the order. There is no
averaging, no last-writer-wins, no silent merge: a genuine collision is a contradiction, and the order admits
no configuration that holds both values at once.

The honest floor is a no-go --- a collision admits no continuation and stops the kernel --- the collision ours
to cause, the wall that answers it the world's. The reading is that
conflicting unordered writes are rejected, not merged. The machine does not paper over a contradiction; faced
with two incompatible values for one slot at one moment, it halts, because no admissible refinement contains
both --- a wall against inconsistency, enforced by stopping.

\section{The Prover--Verifier Effect: a law is the surviving fixed point}
\label{se:prover-verifier}

A physical law can be cast as the survivor of an interactive proof. Modern proof theory casts verification as a
dialogue: a prover proposes, a verifier challenges, and a claim is accepted only if it answers every challenge
the verifier can raise. Cast the space of admissible models as the prover and the record as the verifier, and a
law is what survives every round --- the unique fixed point of the interaction, the model no admissible
challenge can break --- the challenges ours to pose, what survives them the world's. It is not posited and then defended; it is whatever is left standing when every objection
the record can mount has been answered. The surprising strength of the picture is that the verifier need not be
as clever as the prover: by posing challenges it draws at random, a modest verifier can catch a cheating
prover with overwhelming probability --- accepting a true claim and rejecting a false one without ever
reconstructing the whole proof, conviction bought by interrogation rather than by exhaustive check.

The honest floor is a finite count obligation --- the law is the fixed point that survives the challenge
rounds, a survivor counted against every admissible objection --- the interrogation ours, the survivor the
world's. The reading is that a law is what survives
interrogation. It is established by survival rather than by assertion: the law is the fixed point of prover and
verifier, the one model the record cannot refute, left standing when the challenges run out.

\section{The Chaitin Effect: a record with no extractable law}
\label{se:chaitin}

Some records are precise and yet carry no law to extract. Chaitin's halting probability $\Omega$ is the
standing example: a real number, perfectly definite --- the probability that a program assembled at random
halts --- whose digits are algorithmically random, meaning no program shorter than the digits themselves can
produce them. This is the language of Kolmogorov complexity: the complexity of a string is the length of the
shortest program that outputs it, and a string is random when that shortest program is no shorter than the
string itself --- there is nothing to say but the string. Most strings are random in this sense, incompressible
because there are far more strings than short programs to generate them, and $\Omega$ is the exemplar whose
incompressibility can be proved. Its leading digits secretly encode the answers to the halting problem, which is exactly why no
rule can shorten them: to compress $\Omega$ would be to decide what cannot be decided --- the shortening ours
to attempt, the digits that refuse it the world's. Two records can agree on
every digit of a prefix and disagree on the very next, with no predictor that could have called it; the record
is incompressible, its next digit irreducible to any rule the past licenses.

The honest floor is name-only --- the irreducibility is labelled, the algorithmic randomness written as the
mathematics it is, with the finite model claiming only the name --- the label ours, the lawlessness the
world's. The reading is that a precise record can still
hold no law. Precision is not lawfulness: $\Omega$ is exact to every digit and compressible to none, and a
record can be perfectly definite and yet carry nothing a rule could continue --- the limit of extraction,
named.

\section{The Cantor--G\"odel--Cohen Effect: a hypothesis the refinement cannot decide}
\label{se:cantor-godel-cohen}

Some questions no faithful refinement settles. Begin with Cantor. He proved that the continuum --- the real
numbers --- is strictly larger than the integers, by the diagonal argument: take any supposed enumeration of
the reals between $0$ and $1$, one per row, and build a new real differing from the first in its first digit,
the second in its second, the $n$th in its $n$th; the result lies on no row of the list, so no list exhausts
the reals, and $|\mathbb{R}| > |\mathbb{N}|$. There are at least two sizes of infinity. Cantor then asked
whether any infinity lies strictly between them --- the continuum hypothesis, the claim that there is none, that
the reals are the very next size up from the integers.

That a precise question might have no answer the axioms determine was itself a discovery, and G\"odel made it
first. His incompleteness theorems, of 1931, showed that any consistent formal system rich enough to express
arithmetic contains statements that are true but unprovable within it --- the system cannot settle every
question it can pose, and adding axioms to settle one only opens others. The continuum hypothesis turned out to
be a question set theory cannot settle --- not true-but-unprovable, as G\"odel's own examples were, but
undecidable in both directions --- and its undecidability, established in two halves, is the standing instance
of that general wall.

The hypothesis resisted proof for sixty years, and the reason came in two parts. G\"odel showed, in 1940, that
the continuum hypothesis cannot be refuted from the standard axioms of set theory: assume it, and no
contradiction follows. Cohen showed, in 1963, by the method he invented and called forcing, that it cannot be
proved from them either: one can build a model of the axioms in which the hypothesis fails. Forcing adjoins to a
model new sets approached by finite conditions --- a generic object the model itself cannot name, built up by
finite approximations --- and tunes them to make the hypothesis false while keeping every axiom true. Between
the two results the hypothesis is independent: neither true nor false relative to the axioms, but a free choice
of completion. One may extend the axioms to make it hold or to make it fail, and no faithful refinement of the
record decides between them, while an unfaithful one --- which silently inflates the count --- is detected ---
the axioms ours to extend, the independence the world's.

The honest floor is a no-go --- the continuum hypothesis is independent, undecidable by any faithful refinement,
and its resolution is a chosen completion rather than a derived fact --- the completion ours, the
undecidability the world's. The reading is that completion is a
choice, not a discovery. Where the record is genuinely undecided, settling it is an act of extension, not of
measurement; the continuum hypothesis is the standing case --- a question the apparatus cannot decide by
refining, only by choosing how to complete.

\section{The von Neumann Effect, I: the machine-epsilon floor}
\label{se:von-neumann-floor}

A finite machine has a smallest distinguishable difference, and below it computation cannot go. Every number
the machine holds is rounded to a finite precision, and the gap between $1$ and the next representable number is
the machine epsilon, $\epsilon_{\text{mach}}$ --- the floor of the machine's arithmetic. Once rounding error
dominates, once a step needs more precision than $\epsilon_{\text{mach}}$ allows, further operations do not
refine
the answer; they only stir the noise. How fast that error grows is measured by the condition number $\kappa$
of the operation: a well-conditioned step, $\kappa$ near $1$, passes its input's error through roughly
unchanged, while an ill-conditioned step, $\kappa$ large, amplifies a tiny input error into a large output one.
Solve a nearly singular system, or subtract two nearly equal large numbers, and $\kappa$ is enormous --- the
few digits of agreement are eaten by the rounding, and the answer is noise dressed as precision. An
ill-conditioned step, one that demands more precision than is accessible, is inadmissible: it cannot be carried
out faithfully --- the demand ours, the floor the world's.

The honest floor is a no-go --- when the required precision exceeds the accessible precision, the step is
blocked, and the machine-epsilon floor is a wall --- the step ours to pose, what the rounding eats the
world's. The reading is that finite precision is a floor with teeth. A
computation that needs more precision than $\epsilon_{\text{mach}}$ allows cannot be carried out, and the
condition number says in advance which computations the floor will swallow --- the ill-conditioned demand a
precision the machine does not have.

\section{The von Neumann Effect, II: stability as survivorship}
\label{se:von-neumann-stability}

Which computations survive finite precision is subtler than conditioning alone, and Trefethen's answer is the
pseudospectrum. For a well-behaved, normal operator the eigenvalues --- the spectrum --- tell how perturbations
grow: a stable spectrum, a stable computation. But many operators are non-normal, and for them the spectrum
lies. An operator can have every eigenvalue safely inside the stable region and still amplify perturbations
enormously over a long transient, because its eigenvectors are nearly parallel and a small input projects onto
a huge combination of them. The pseudospectrum captures what the spectrum hides: it is the set of values where
the resolvent, the inverse $(zI - A)^{-1}$, is merely large, not only where it is infinite as at the spectrum
proper. A pseudospectrum bulging into the unstable region warns of growth the eigenvalues conceal --- the
promise ours to take from the spectrum, the survivorship the world's.

The honest floor is a no-go --- a computation whose pseudospectrum reaches into the unstable region does not
survive finite precision, however safe its spectrum looks --- the spectrum ours to read, what survives the
world's. The reading is that stability is survivorship, not
spectrum. What lives under rounding is what the pseudospectra permit; the spectrum can promise a stability the
operator does not have, and only the resolvent, read where it is large, says which computations the floor lets
through --- stability the name for surviving the perturbations finite precision injects.

\section{The Newton--Cooley--Tukey Effect: recursion divides and conquers}
\label{se:newton-cooley-tukey}

The machine's reach is extended by recursion. Two classics show the pattern. The fast Fourier transform factors
a transform of size $N$ into two of size $N/2$, and those again, computing in $O(N \log N)$ what would naively
take $O(N^2)$ --- the recursion halving the work at each level, a saving that turned the transform from a
curiosity into the workhorse of signal processing. And Newton's iteration refines a root by repeated
correction,
\begin{equation}
  x_{n+1} = x_n - \frac{f(x_n)}{f'(x_n)},
\end{equation}
each step sliding down the tangent line to where it crosses zero, and each step roughly squaring the accuracy
--- a handful of corrections from a rough guess to machine precision. Both are divide-and-conquer made into
admissible refinement: split the problem, solve the parts, combine --- the splitting ours, the reach it buys
the world's.

The honest floor is a finite count obligation --- the recursion is a count of admissible steps, $O(N\log N)$
for the transform and a handful of corrections for the root --- the factoring ours to run, the count of steps
the world's. The reading is that recursion is refinement that
divides and conquers. The machine extends its reach not by brute force but by factoring a problem into smaller
admissible steps, and the fast transform and Newton's method are the same move twice --- divide, refine the
parts, combine, each step a legitimate refinement of the record.

\bigskip

With recursion, the machine's one power and four walls stand clear. It can
replay anything it can serialize --- that is the universal power (Turing). Against it stand the four walls. It
must stop where the order admits no successor and where writes collide (Halt and Jupyter): the wall of
continuation. It cannot compress what is algorithmically random (Chaitin): the wall of extraction. It cannot
decide the independent (Cantor--G\"odel--Cohen): the wall of decision. And it cannot stabilize the
ill-conditioned or carry the non-normal through finite precision (von Neumann, floor and survivorship both):
the wall of stability. Where a problem factors, recursion carries it (Newton--Cooley--Tukey), and a law,
between the walls, is what survives interrogation (Prover--Verifier) --- but the four walls do not factor away
--- the program still ours to write, the walls still the world's.
Run as a computer, the apparatus gains universality and meets exactly these four limits. Part VI turns next to
the scaling of systems, and the apparatus toward its last examination.

\bigskip
\begin{center}
\rule{0.4\textwidth}{0.4pt}
\end{center}
\bigskip

\noindent The case comes to its last examination, and it is the deepest: it runs its own apparatus as a machine,
and asks what a finite reckoning can do and where it must stop. The instrument that has read the whole world is
set, at the end, to compute --- and a computer has exactly one power and exactly four walls.

The power is universal. Anything the instrument can write down as a sequence it can replay: flatten a record to
a line, hand the line to a single reader that takes one step per mark, and every record ever kept is within its
reach. Against that one power stand four walls, and each is a wall of the same honest kind the book has met
since the friction threshold --- a place the machine runs up to and cannot pass, named plainly and not
pretended away. It cannot always continue: where the order admits no next step, the reckoning halts --- complete
not by triumph but because there is nowhere left to go. It cannot always compress: some records are exact to
every digit and carry no law shorter than themselves, incompressible, their next mark irreducible to any rule
the past licenses. There is even a number that is the plainest case of it --- the odds that a program thrown
together at random ever halts, exact to every digit and shorter than none, its leading digits secretly holding
which programs stop, so that to compress it would be to decide what cannot be decided. It cannot always decide: some questions no faithful refinement settles, true or false only
once a choice is made that the record itself does not force. And it cannot always stay stable: below the
machine's own smallest step, a reckoning that needs more precision than it has does not sharpen the answer but
only stirs the noise. The program is ours to write; the four walls it runs into are the world's.

The third wall is the one the whole book has been walking toward, and here at last it is named. It was shown
long ago that the points on a line outnumber the whole numbers --- no list, however long, can name them all, for
one can always write down a point missing from every row of it. Ask, then, whether between those two sizes of
infinity there lies any third in between --- and the question has no answer the axioms decide. It was proved, in
one half, that assuming it breeds no contradiction, and, decades later, in the other, that a world can be built
obeying every axiom in which it fails; one can have a model where it holds and a model where it does not, and no
faithful refinement of the record says which is the true one. The way that second world
is built has a name, given by the man who found it: \emph{forcing}. To force is to adjoin to a model a new
object the model cannot name --- a generic, reached not all at once but built up out of finite conditions,
finite approximations closing in on a thing the finite record can never itself complete. And here the book's own
method stands revealed. This is what the instrument has been doing since the first mark: laying down finite
conditions --- counts, one at a time --- and reaching, through them, for a world it approaches and never
finishes, a generic beyond the last mark it will ever make. The book has been forcing all along, and did not say
so until now. Where the record is genuinely undecided it does not measure the answer; it chooses a completion
--- the axioms ours to extend, the undecidability the world's.

And the \emph{charge} is that reckoning, run to its limit. The count the case will collect on is a finite
reckoning --- a machine's worth of steps, and no more --- and it stands among the four walls, honest about each:
it does not claim to continue where the order stops, nor to compress what carries no law, nor to decide what no
refinement settles, nor to hold steady below its own smallest step. What it owes, it owes within the walls,
computed as far as the machine faithfully reaches and marked, at each wall, as a thing it met and did not cross.
A charge honest about its own limits is worth more than one that pretends to none, and the instrument has spent
the whole book earning exactly that honesty --- the charge inherits it. The trace is the reckoning entire, walls
and all: the count carried as far as it could faithfully go, and the four places it stopped, each named. What
comes to be owed is owed by a finite machine that knows precisely what it cannot do.

And with this the instrument has examined itself to the end. It has asked what makes a reading a measurement,
how finely it can be made, whether it can be carried forward, how its order and its geometry arise, what algebra
its records obey, and now what machine runs on them and where that machine must stop. There is nothing left of
the apparatus it has not turned upon itself. The record is complete. The three readings the case was built on
are gathered, corrected, floored, and closed under one algebra; they run on one machine and meet its four walls;
and they can, at last, be brought together. Everything the case needed is in hand --- save the one thing no
examination could supply: the meeting itself, the three made one, the number they must agree on.

So Part VI closes, and with it the apparatus's long turn upon itself. Every floor has been named --- the count,
the shadow, the wall, the label --- every effect fenced at exactly the one it stood on; every wall has been met
and marked; and the instrument, having read the whole world and then itself, has nothing left to examine. The
record is complete at last. For the first time in thirty chapters, the court can sit. What remains is not another
reading but the reckoning of them all --- the three readings brought to one, and that one number carried, at
last, to the bench. The book has a single chapter more, and it is the trial. The court is ready to sit; the case
is ready to be heard.
